\section{Limitations}
\label{sec:limitations}

This survey has several limitations that qualify our conclusions.

\textbf{Recency bias.} The field evolves rapidly; new systems appear
weekly. Our survey covers publications and releases through February
2026. Systems published after this date are not included.

\textbf{Selection bias.} We identify systems primarily through
English-language academic databases, GitHub, and practitioner blogs.
Systems published in non-English venues or deployed in closed-source
products may be underrepresented.

\textbf{Classification subjectivity.} Mapping systems to the CoALA
taxonomy, RAG categories, and lifecycle operations requires
interpretive judgment. We mitigate this through explicit definitions
(\Cref{sec:taxonomy}) and consistent application, but reasonable
disagreements are possible.

\textbf{Benchmark comparability.} Not all systems report results on
the same benchmarks. LoCoMo, LongMemEval, and DMR scores are not
directly comparable across papers due to variations in evaluation
setup, model choice, and preprocessing.

\textbf{Reproducibility.} All baseline experiments
(\S\ref{sec:baseline-replication}, \S\ref{sec:mab-replication}) use
Gemini-2.0-Flash-001 via OpenRouter (accessed February 2026;
temperature~0, max\_tokens~512) as the generator and GPT-4o-mini via
OpenRouter (temperature~0, max\_tokens~10) as the judge.  Embedding
model: \texttt{all-MiniLM-L6-v2} via ChromaDB.  Total inference cost
for all three benchmarks: approximately \$12.  Code and raw results are
available at~\cite{alaya2026}.  Model
behavior may change with API updates; our results reflect model versions
available in February 2026.

\textbf{No system-level evaluation.} This paper surveys and classifies
memory architectures but does not evaluate any specific system's
end-to-end performance. Comparative evaluation of lifecycle operations
(forgetting, consolidation, transformation) across systems remains
future work.

\textbf{Evolving benchmark landscape.} Since our evaluation cutoff,
several benchmarks have appeared that address gaps we identify in
\S\ref{sec:benchmark-landscape} and \S\ref{sec:open-problems}.
AMA-Bench~\cite{amabench2026} evaluates long-horizon agentic memory
with trajectory processing across extended interaction sequences, going
beyond the static-history paradigm of LoCoMo and LongMemEval. The
Agent Memory Benchmark (AMB) from Hindsight/Vectorize extends
LoCoMo and LongMemEval to agentic task settings where memory must
inform action selection, not just answer retrieval questions. These
benchmarks, together with MemoryAgentBench~\cite{hu2025mab}, represent
the field's shift toward evaluating memory as a \emph{process} rather
than a \emph{store}. Future work should evaluate systems on this
expanded benchmark suite.

\textbf{Temporal reasoning remains catastrophic.} Our baseline
experiments (\S\ref{sec:baseline-replication}) confirm that temporal
reasoning accuracy ranges from 6.5\% to 29.3\% across all systems and
benchmarks tested---well below random baselines on binary questions.
TReMu~\cite{ge2025tremu} corroborates this finding. No surveyed
production system addresses temporal reasoning as a first-class
capability, and this remains the single largest accuracy gap in the
field.

\textbf{Memory staleness detection.} No surveyed system implements
mechanisms to detect when high-strength memories become factually
outdated. A memory encoded with high storage strength (e.g., ``user
works at Company~X'') persists indefinitely even after the underlying
fact changes. Staleness detection---identifying memories whose
real-world referents have shifted---requires temporal grounding and
external validation that current architectures do not provide.

\textbf{Multi-agent memory coordination.} The market is moving toward
multi-agent architectures (CrewAI, AutoGen, Agency Swarm), but most
memory systems---including all those evaluated in this survey---assume a
single-agent context. How multiple agents should share, partition, or
synchronize a common memory store, and how forgetting and consolidation
policies should interact across agents with different roles and
information needs, are open problems with no published solutions.
